Innerhalb der Tabelle bedeutet \(n=\lambda\cdot\mu\).
\begin{tabproofwide}
  \proofstepwidestarbreak[1]{\FiniteSemigroupAlg{A}{\star}}{\rA}
  \proofstepwidestarbreak[2]{a\in A}{\rA}
  \proofstepwidestarbreak[1]{\SemigroupAlg{A}{\star}}{\FormulaRefAuto{FiniteSemigroupBaseAxiom}[ax]{1}}
  \proofstepwidestarbreak[1,2]{\exists\mu\,\exists\lambda\,\operatorname{IP}_{A,\star}(a;\mu,\lambda)}{\FormulaRefAuto{FiniteSemigroupIndexPeriodExistence}[thm]{1,2}}
  \proofstepwidestarbreak[5]{\exists\lambda\,\operatorname{IP}_{A,\star}(a;\mu,\lambda)}{\rA}
  \proofstepwidestarbreak[6]{\operatorname{IP}_{A,\star}(a;\mu,\lambda)}{\rA}
  \proofstepwidestarbreak[1,2,6]{\mu\in\mathbb N\land\mu\neq0\land\lambda\in\mathbb N\land\lambda\neq0}{\FormulaRefAuto{FiniteSemigroupIndexPeriodPositiveData}[ax]{3,2,6}}
  \proofstepwidestarbreak[1,2,6]{a^{\mu+\lambda}=a^\mu}{\FormulaRefAuto{FiniteSemigroupIndexPeriodEquality}[ax]{3,2,6}}
  \proofstepwidestarbreak[1,2,6]{n\in\mathbb N\land n\neq0\land\mu\leq n}{\FormulaRefAuto{FiniteSemigroupIdempotentExponentData}[thm]{\rAI{\rAEa{7},\rAEa{\rAEb{\rAEb{7}}}},\rAEa{\rAEb{7}},\rAEb{\rAEb{\rAEb{7}}}}}
  \proofstepwidestarbreak[1,2,6]{a^{n+n}=a^n}{\FormulaRefAuto{FiniteSemigroupMultiplePeriodShift}[thm]{3,2,\rAI{\rAEa{7},\rAI{\rAEa{\rAEb{\rAEb{7}}},\rAI{\rAEa{9},\rAEa{7}}}},\rAEa{\rAEb{7}},\rAEb{\rAEb{\rAEb{7}}},8,\rAEb{\rAEb{9}}}}
  \proofstepwidestarbreak[1,2,6]{a^{n+n}=a^n\star a^n}{\FormulaRefAuto{SemigroupPowerAdditionLaw}[thm]{3,2,\rAEa{9},\rAEa{9},\rAEa{\rAEb{9}},\rAEa{\rAEb{9}}}}
  \proofstepwidestarbreak[1,2,6]{a^n\in A}{\FormulaRefAuto{SemigroupPowerClosure}[thm]{3,2,\rAEa{9},\rAEa{\rAEb{9}}}}
  \proofstepwidestarbreak[1,2,6]{a^n\star a^n=a^n}{\rChain{\FormulaRefAuto{a=b\vdash b=a}[thm]{11},10}}
  \proofstepwidestarbreak[1,2,6]{a^n\in E_{A,\star}}{\FormulaRefAuto{SemigroupIdempotentMembership}[thm]{3,\rAI{12,13}}}
  \proofstepwidestarbreak[1,2,6]{\exists n\in\mathbb N\,(n\neq0\land a^n\in E_{A,\star})}{\rEI{\rAI{\rAEa{9},\rAI{\rAEa{\rAEb{9}},14}}}}
  \proofstepwidestarbreak[1,2,5]{\exists n\in\mathbb N\,(n\neq0\land a^n\in E_{A,\star})}{\rEE{5,6,15}}
  \proofstepwidestarbreak[1,2]{\exists n\in\mathbb N\,(n\neq0\land a^n\in E_{A,\star})}{\rEE{4,5,16}}
\end{tabproofwide}
